% Diagramma di corpo libero di una lampada appesa a due funi, ciascuna a 60 gradi
% dalla verticale: le due tensioni F_T lungo le funi e il peso F_p. Per la fune
% di destra sono tratteggiate le componenti F_T cos 60 (verticale) e
% F_T sin 60 (orizzontale). In equilibrio F_T = F_p, e le tre frecce hanno la
% stessa lunghezza.
\begin{tikzpicture}[>=Stealth, font=\small, scale=1.3]
  \coordinate (K) at (0,0);
  % soffitto e funi
  \fill[richiamo!25] (-4.0,2.0) rectangle (4.0,2.15);
  \draw[thick] (-4.0,2.0) -- (4.0,2.0);
  \draw[thick] (K) -- (30:4.0);
  \draw[thick] (K) -- (150:4.0);
  % lampada
  \filldraw[fill=esempio!35, draw=esempio, thick] (-0.2,-0.1) -- (0.2,-0.1) -- (0.5,-0.6) -- (-0.5,-0.6) -- cycle;
  \draw[thick] (K) -- (0,-0.1);
  \fill (K) circle (1.4pt);
  % verticale di riferimento e angolo di 60 gradi
  \draw[dotted] (K) -- (0,1.9);
  \draw (0,0.7) arc (90:30:0.7);
  \node at (60:0.95) {\ang{60}};
  \draw (0,0.7) arc (90:150:0.7);
  \node at (120:0.95) {\ang{60}};
  % forze
  \draw[->, formula, very thick] (K) -- (30:2.0) coordinate (T) node[below right, xshift=2pt] {$\vec F_T$};
  \draw[->, formula, very thick] (K) -- (150:2.0) node[below left] {$\vec F_T$};
  \draw[->, attenzione, very thick] (K) -- (0,-2.0) node[right] {$\vec F_p$};
  % componenti della tensione di destra
  \draw[->, formula!70, thick, dashed] (K) -- (T |- K) coordinate (Tx);
  \node[formula, right] at (Tx) {$F_T\sin\ang{60}$};
  \draw[->, formula!70, thick, dashed] (K) -- (K |- T) coordinate (Ty);
  \node[formula, above] at (Ty) {$F_T\cos\ang{60}$};
  \draw[dotted] (Tx) -- (T) -- (Ty);
\end{tikzpicture}
